\documentclass[11pt,a4paper]{article}
\usepackage[utf8]{inputenc}
\usepackage[french]{babel}
\usepackage[T1]{fontenc}
\usepackage{lmodern}
\usepackage[margin=1.5cm]{geometry}

\begin{document}

% En-tête
\begin{center}
    {\Huge \textbf{Emma Dubois}} \\[0.2cm]
    {\Large \textbf{Pentester / Ethical Hacker}} \\[0.1cm]
    emma.dubois@email.com \ | \ 06 95 57 43 62
\end{center}

\vspace{0.3cm}

% Résumé / Profil
\section*{Profil Professionnel}
\noindent
Ingénieur en cybersécurité passionné par la protection des infrastructures critiques et la traque des menaces avancées (Threat Hunting). Doté d'une solide expertise technique en test d'intrusion et en réponse à incident, je mets un point d'honneur à sensibiliser les équipes de développement aux enjeux du DevSecOps. Mon objectif est de garantir l'intégrité et la confidentialité des données sensibles face aux cyberattaques.

\vspace{0.3cm}

% Compétences
\section*{Compétences Techniques}
\noindent
\textbf{Wireshark} $\bullet$ \textbf{DevSecOps} $\bullet$ \textbf{Bash} $\bullet$ \textbf{Cryptographie} $\bullet$ \textbf{ISO 27001} $\bullet$ \textbf{Splunk (SIEM)} $\bullet$ \textbf{Forensic} $\bullet$ \textbf{Pentest}

\vspace{0.3cm}

% Expériences
\section*{Expériences Professionnelles}
\noindent \textbf{Pentester / Ethical Hacker / Confirmé} --- \textit{Airbus CyberSecurity} \hfill \textbf{05/2025 à Aujourd'hui} \\
Reverse engineering de souches virales inédites pour extraire les indicateurs de compromission (IOC) et mettre à jour les bases de signatures. Déploiement et configuration avancée d'une solution EDR (Endpoint Detection and Response) sur un parc de plus de 5 000 postes de travail. Gestion des incidents de sécurité de criticité majeure (P1) en coordination avec l'équipe de réponse à incident (CERT) et le comité de direction. Implémentation d'une approche DevSecOps en intégrant des outils d'analyse statique et dynamique (SAST/DAST) directement dans les pipelines CI/CD. \\[0.3cm]
\noindent \textbf{Pentester / Ethical Hacker} --- \textit{Sogeti} \hfill \textbf{02/2023 à 03/2025} \\
Accompagnement à la mise en conformité ISO 27001, incluant la rédaction des politiques de sécurité (PSSI) et la cartographie des risques. Animation de campagnes de sensibilisation à la cybersécurité (phishing simulé) ayant permis de réduire le taux de clics dangereux de 40\%. Analyse en temps réel des alertes de sécurité remontées par le SIEM (Splunk), permettant la détection et le blocage d'attaques par ransomware. Conception d'architectures réseaux sécurisées (Zero Trust) pour les environnements Cloud AWS, avec restriction drastique des privilèges IAM. \\[0.3cm]
\noindent \textbf{Stage — Assistant Pentester / Ethical Hacker} --- \textit{Orange Cyberdefense} \hfill \textbf{06/2021 à 02/2023} \\
Investigation numérique (Forensic) suite à une compromission de serveur, identification de la porte dérobée (backdoor) et restauration de l'intégrité du système. Réalisation de plus de 20 tests d'intrusion (pentests) sur des applications web et des infrastructures internes, avec rédaction de rapports d'audit détaillés. Création de scripts d'automatisation en Python pour accélérer la collecte d'artefacts malveillants et l'analyse de logs réseau. \\[0.3cm]


\vspace{0.3cm}

% Formations
\section*{Formation \& Diplômes}
\noindent \textbf{Diplôme d'Ingénieur en Cybersécurité} \hfill \textbf{2020 - 2023} \\
\textit{ESIEA} \\[0.3cm]
\noindent \textbf{Classes Préparatoires (CPGE) en Filière Scientifique} \hfill \textbf{2018 - 2020} \\
\textit{Lycée Scientifique} \\[0.3cm]


\end{document}